% Generated draft: requires mathematical and source review.
\documentclass[11pt]{article}
\usepackage[margin=25mm]{geometry}
\usepackage{fontspec}
\setmainfont{Latin Modern Roman}
\usepackage{amsmath,amssymb,mathtools,unicode-math}
\setmathfont{Latin Modern Math}
\usepackage{enumitem,needspace,xurl,hyperref}
\setlength{\parindent}{0pt}
\setlength{\parskip}{5pt}
\setlength{\emergencystretch}{3em}
\newcommand{\N}{\mathbb{N}}
\newcommand{\Z}{\mathbb{Z}}
\newcommand{\Q}{\mathbb{Q}}
\newcommand{\R}{\mathbb{R}}
\newcommand{\C}{\mathbb{C}}
\newcommand{\F}{\mathbb{F}}
\newcommand{\A}{\mathbb{A}}
\newcommand{\PP}{\mathbb{P}}
\newcommand{\E}{\mathbb{E}}
\newcommand{\eps}{\varepsilon}
\newcommand{\dd}{\,\mathrm{d}}
\begin{document}
\section*{UnsolvedMath ID: 701166}
\textbf{Source problem code:} catalogue entry
\subsection{Definitions and mathematical statement}
Let \(n\) be a positive integer, let \(F\) be a free group of any rank, and let \(c:F\to\{1,\ldots,n\}\) be an \(n\)-colouring. A subset \(S\subseteq F\) is monochromatic if there is an \(i\in\{1,\ldots,n\}\) such that \(c(s)=i\) for every \(s\in S\).

The source does not define symmetry with respect to a group element. Under the usual point-reflection interpretation, symmetry with respect to \(g\in F\) means invariance under the involution \(s\mapsto gs^{-1}g\), or equivalently
\[
S=gS^{-1}g,\qquad S^{-1}=\{s^{-1}:s\in S\}.
\]
Under this interpretation, and without adding a size requirement absent from the record, the supplied question has the form
\[
\begin{gathered}
\text{For every positive integer }n,\text{ every free group }F\text{ of any rank,}\\
\text{and every }c:F\to\{1,\ldots,n\},\text{ do there exist }g\in F,\ S\subseteq F,\ i\in\{1,\ldots,n\}\\
\text{such that }S=gS^{-1}g\quad\text{and}\quad (\forall s\in S)\ c(s)=i\ ?
\end{gathered}
\]
The source supplies no requirement that \(S\) be nonempty, infinite, or of any prescribed cardinality. Its intended size condition therefore remains unresolved; the display records only the literal size-unrestricted reading under the stated symmetry convention. In particular, an infinitude requirement must not be silently inserted. If such a requirement was intended, the permitted ranks would also need clarification. The source states that the assertion is true for \(n=2\), but the precise size and symmetry conventions governing that assertion are not supplied.
\subsection{Short English statement}
The question asks whether every finite colouring of a free group of any rank has a monochromatic subset symmetric about some group element. The source states a positive result for two colours. The meaning of symmetry and the required size of the subset need clarification before the intended problem can be stated precisely.
\subsection{Sources}
\begin{itemize}
\item Local UnsolvedMath registry: supplied complete source record, permanent ID \(701166\), exact source problem code \texttt{catalogue entry}; upstream revision \texttt{d371ee917ac9647fbb06cd6fc529deec6f3dda82}.
\item Supplied reference: \emph{The Kourovka Notebook (Unsolved Problems in Group Theory)}, problem 15.78, 15th issue (2002), retained in the supplied 21st-issue edition (2026), arXiv:1401.0300v48. The collection metadata names E. I. Khukhro and V. D. Mazurov as editors. No external reference was independently consulted.
\item Attribution: the background and proposer field name I. V. Protasov; the statement itself begins with the parenthetical attribution R. I. Grigorchuk. Their respective roles are not explained in the supplied record.
\item The registry labels the problem \texttt{open} and supplies the assertion that the two-colour case is true. These are supplied source claims, not independently verified mathematical or literature assessments.
\end{itemize}
\end{document}
